\section{Optimal element size}
\label{sec:optimal-size}

\subsection{A flexible toy primitive}

Fix a total payload size $P$ bytes and a state width $W$ bytes, and suppose
the element size $D$ can be freely chosen (i.e.\ the compression function's
block width can be made to match $D$ exactly, as in our toy 32-byte
primitive, Section~\ref{sec:examples}). $D$ is not a free real number,
however: $\mathrm{MID}(a)$ re-injects the entire $W$-byte checkpoint into
one $D$-byte block via an injective encoding
(Definition~\ref{def:padb}), and no injective encoding can embed a wider
value into a narrower one, so
\begin{equation}
D \ge W
\end{equation}
is a hard feasibility floor, not a property of any particular padding
scheme -- it holds regardless of whether the encoding used has a spare
terminator byte (the $b>w$ scheme) or none (the $b=w$ identity encoding).
The construction also requires $N = 2^L$ for an integer $L$
(Section~\ref{sec:construction}; non-power-of-two $N$ is not yet handled --
Section~\ref{sec:non-power-of-two}), so the genuine free variable is the
integer $L$, with $D = P/2^L$ determined by it, not $D$ itself. The problem
is therefore the constrained optimization
\begin{equation}
\min_{L \in \mathbb{Z}}\ \mathrm{size}(L) = 2D + (L-1)W = \frac{2P}{2^L} + (L-1)W
\qquad \text{subject to } D = P/2^L \ge W,
\end{equation}
two raw elements plus $L-1$ state reveals, minimized only over the integer
$L$ values whose resulting $D$ actually satisfies the floor.

Ignoring the constraint first, to see where it binds: treating $L$ as
continuous and minimizing over (continuous) $D = P/2^L$ gives the
stationarity condition
\begin{equation}
\frac{d}{dD}\,\mathrm{size}(D) = 2 - \frac{W}{D \ln 2} = 0
\quad\Longrightarrow\quad
D^* = \frac{W}{2\ln 2}.
\end{equation}
For $W = 32$ bytes, $D^* \approx 23.08$ bytes -- strictly \emph{less} than
$W=32$, i.e.\ the unconstrained stationary point lies entirely outside the
feasible region $D \ge W$. Since $\mathrm{size}(D)$ is convex and its
derivative is already positive at $D=W$ (the constraint boundary), the
constrained optimum is not this stationary point at all, but simply the
\emph{smallest feasible} $D$: $D = W$ exactly, using the identity encoding.
Every power of two below $W$ is infeasible outright -- not merely
suboptimal -- because the checkpoint cannot physically fit.

\begin{table}[H]
\centering
\begin{tabular}{@{}rrrrl@{}}
\toprule
$L$ & $N = 2^L$ & $D = P/N$ (bytes) & Proof size (bytes) & Feasible? \\
\midrule
21 &   2{,}097{,}152 & 512.00 & 1664.00 & yes \\
22 &   4{,}194{,}304 & 256.00 & 1184.00 & yes \\
23 &   8{,}388{,}608 & 128.00 &  960.00 & yes \\
24 &  16{,}777{,}216 &  64.00 &  864.00 & yes \\
\textbf{25} &  \textbf{33{,}554{,}432} &  \textbf{32.00} &  \textbf{832.00} & \textbf{yes ($D=W$)} \\
26 &  67{,}108{,}864 &  16.00 &  (832.00) & \textbf{no: $D<W$} \\
27 & 134{,}217{,}728 &   8.00 &  (848.00) & \textbf{no: $D<W$} \\
28 & 268{,}435{,}456 &   4.00 &  (872.00) & \textbf{no: $D<W$} \\
\bottomrule
\end{tabular}
\caption{Sweep at fixed payload $P = 2^{30}$ bytes (1024 MiB), $W=32$
bytes. Rows with $D<W=32$ (bracketed sizes) are not achievable FlatMerkle
configurations at all -- a 32-byte checkpoint cannot be injectively
embedded in a smaller block, under any encoding -- not merely worse than
the optimum; their $\mathrm{size}(L)$ values are shown only to illustrate
that the unconstrained continuous formula would (wrongly) suggest they
help. The true, unique discrete minimum is $L=25$ ($D=32=W$, using the
identity checkpoint encoding of Definition~\ref{def:padb}), 832.00 B --
not a tie with any smaller $D$, and not the non-power-of-two, sub-$W$
point $D\approx23.08$ the unconstrained continuous relaxation alone would
suggest. A deployment that instead needs the $b>w$ terminator-byte
encoding (e.g.\ to keep one spare byte for a domain-separation tag,
Remark~\ref{rem:domain-sep-bit}) faces the stricter floor $D \ge W+1=33$
bytes; since $D$ must itself be a power of two here ($D=P/2^L$ with both
$P$ and $N$ powers of two), the smallest such $D$ is $64$, i.e.\ $L=24$,
864.00 B.}
\label{tab:optimal-toy}
\end{table}

\subsection{Real SHA-256: the block size is fixed at 64 bytes}
\label{sec:optimal-real}

The toy analysis assumes $D$ can be chosen freely, but SHA-256's block
width is fixed at $b=512$ bits (64 bytes): every item of the structured
file, data or $\mathrm{MID}$, occupies exactly one 64-byte block, and
Section~\ref{sec:construction} does not pack several application records
into one element. With a payload of $P$ bytes split into 64-byte elements,
$N = P/64$, and since each proof reveal (a checkpoint) carries only its
32-byte state, a single-element proof costs
\begin{equation}
\mathrm{size}(P) = 2 \times 64 + \big(\log_2(P/64) - 1\big) \times 32
= 96 + 32\,\log_2(P/64) \text{ bytes}.
\end{equation}
This was verified directly against the real-SHA-256 implementation: for
$N=32$ real 64-byte elements ($P = 2048$ bytes), the formula predicts
$96 + 32\log_2(32) = 96 + 160 = 256$ bytes, which matches the measured
proof size exactly.

Unlike the toy case, where the two coincide at $D=W$, here FlatMerkle's
proof is \emph{larger} than a Merkle proof over the same 64-byte elements:
a Merkle proof carries one 64-byte leaf and $\log_2 N$ 32-byte siblings,
$64 + 32\log_2(P/64)$ bytes, exactly 32 bytes less than FlatMerkle's, and
it discloses no neighbouring element
(Table~\ref{tab:real-sha256-sizes}). The difference is the second raw
element of FlatMerkle's base pair (Section~\ref{par:pair-disclosure}).

\begin{table}[htbp]
\centering
\begin{tabular}{@{}rrrr@{}}
\toprule
$N$ & $P = 64N$ (bytes) & FlatMerkle proof (bytes) & Merkle proof (bytes) \\
\midrule
4  & 256  & 160 & 128 \\
8  & 512  & 192 & 160 \\
16 & 1024 & 224 & 192 \\
32 & 2048 & 256 & 224 \\
64 & 4096 & 288 & 256 \\
\bottomrule
\end{tabular}
\caption{Real-SHA-256 FlatMerkle proof sizes, measured directly against the
\texttt{loglink\_annex.py} implementation, matching
$96 + 32\log_2(P/64)$ exactly at every scale tested, next to the size of a
Merkle proof over the same 64-byte elements, $64 + 32\log_2 N$ (computed).}
\label{tab:real-sha256-sizes}
\end{table}
